\documentclass[12pt]{article}
% AAU theme adapted from the preamble supplied by the project team.
\usepackage[a4paper,margin=1in]{geometry}
\usepackage{lmodern}
\usepackage[T1]{fontenc}
\usepackage{microtype}
\usepackage{graphicx}
\usepackage{float}
\usepackage{amsmath}
\usepackage[table]{xcolor}
\usepackage{titlesec}
\usepackage{fancyhdr}
\usepackage[most]{tcolorbox}
\PassOptionsToPackage{obeyspaces}{url}
\usepackage[hidelinks]{hyperref}
\usepackage{array}
\usepackage{tabularx}
\usepackage{enumitem}
\usepackage{booktabs}
% Vector UML artwork and multipage test tables.
\usepackage{longtable}
\usepackage{pdflscape}
\definecolor{aauRed}{RGB}{164,22,35}
\definecolor{aauGreen}{RGB}{0,122,88}
\definecolor{darkGray}{RGB}{60,60,60}
\definecolor{softGray}{RGB}{245,245,245}
\pagestyle{fancy}\fancyhf{}\fancyfoot[C]{\thepage}
\renewcommand{\headrulewidth}{0pt}\renewcommand{\footrulewidth}{0pt}
\titleformat{\section}{\normalfont\Large\bfseries\color{aauRed}}{\thesection}{0.6em}{}[{\color{aauRed!65}\titlerule[1.1pt]}]
\titlespacing*{\section}{0pt}{12pt}{8pt}
\titleformat{\subsection}{\normalfont\normalsize\bfseries\color{aauRed!85!black}}{\thesubsection}{0.6em}{}
\titlespacing*{\subsection}{0pt}{10pt}{4pt}
\newtcolorbox{objectivebox}{enhanced,colback=aauRed!5,colframe=aauRed!80!black,boxrule=0.3mm,arc=1mm,left=10pt,right=10pt,top=6pt,bottom=6pt,borderline west={2.5pt}{0pt}{aauRed}}
% Each short use/test case stays on one page, including its heading.
\newtcolorbox{contentbox}[1][]{enhanced,colback=softGray,colframe=black!35,boxrule=0.25mm,arc=2mm,left=9pt,right=9pt,top=7pt,bottom=7pt,before skip=9pt,after skip=9pt,title={\textbf{#1}},fonttitle=\bfseries,coltitle=darkGray,colbacktitle=black!8}
\newcommand{\blankline}{\rule{\linewidth}{0.35pt}}
\newcommand{\insertfigure}[3]{\begin{figure}[H]\centering\fbox{\parbox[c][#2][c]{0.88\textwidth}{\centering\textcolor{black!50}{\Large #1}}}\caption{#3}\end{figure}}
\newcommand{\screen}[3][0.98]{\begin{figure}[H]\centering\includegraphics[width=#1\linewidth]{../assets/#2}\caption{#3}\end{figure}}
% Give wide UML models a dedicated landscape page with readable labels.
\newenvironment{umlpage}{%
\newgeometry{margin=0.65in}\begin{landscape}\pagestyle{empty}%
\AddToHook{shipout/foreground}[arena/umlfooter]{%
\begin{tikzpicture}[remember picture,overlay]
\node[rotate=90,anchor=south,inner sep=0pt] at ([xshift=-23pt]current page.east)
{\normalfont\normalsize\thepage};
\end{tikzpicture}}%
}{\end{landscape}\RemoveFromHook{shipout/foreground}[arena/umlfooter]\pagestyle{fancy}\restoregeometry}
\newcommand{\umlfigure}[2]{\begin{figure}[H]\centering
\includegraphics[width=\linewidth,height=13cm,keepaspectratio]{../assets/uml/#1.pdf}
\caption{#2}\end{figure}}
\providecommand{\reportdate}{10 November 2026}
% Two intentionally sized figures fit together with a heading on an A4 page.
\newcommand{\screenpair}[5][0.30]{%
\begin{figure}[H]\centering
\includegraphics[width=0.96\linewidth,height=#1\textheight,keepaspectratio]{../assets/#2}\caption{#3}
\vspace{8pt}
\includegraphics[width=0.96\linewidth,height=#1\textheight,keepaspectratio]{../assets/#4}\caption{#5}
\end{figure}}
\setlength{\parindent}{0pt}\setlength{\parskip}{5pt}
\setlist[itemize]{leftmargin=1.5em,itemsep=2pt,topsep=3pt}
\setlist[enumerate]{leftmargin=1.7em,itemsep=3pt,topsep=4pt}
\setlength{\emergencystretch}{2em}
\widowpenalty=10000\clubpenalty=10000\displaywidowpenalty=10000
\brokenpenalty=10000
\raggedbottom
\renewcommand{\arraystretch}{1.16}
\setlength{\abovecaptionskip}{7pt}
\setlength{\belowcaptionskip}{4pt}
\setlength{\intextsep}{12pt plus 2pt minus 2pt}
\newcolumntype{Y}{>{\raggedright\arraybackslash}X}
\BeforeBeginEnvironment{tabularx}{\par\noindent}
\AfterEndEnvironment{tabularx}{\par}
\newcommand{\code}[1]{{\small\expandafter\nolinkurl\expandafter{\detokenize{#1}}}}
\newcommand{\command}[1]{\mbox{\small\ttfamily\detokenize{#1}}}
\newcommand{\projecttitle}[2]{%
\hypersetup{pageanchor=false}\begin{titlepage}\centering\vspace*{-0.45cm}
\noindent\begin{minipage}[c]{0.66\textwidth}\raggedright
{\Large\bfseries Al Ain University (AAU)\par}
\end{minipage}\hfill
\begin{minipage}[c]{0.28\textwidth}\raggedleft
% The supplied PNG is unchanged; trim removes only its transparent margins.
\includegraphics[width=3.1cm,trim=1731bp 871bp 1735bp 878bp,clip]{../../aau.png}
\end{minipage}\par
\vspace{0.45cm}
{\scshape\large Object Oriented Analysis and Design\par}
\vspace{0.10cm}
{\normalsize System Analysis \& Design\par}
\vspace{0.10cm}
{\small Course codes: 0103430 \;|\; 0102331\par}
\vspace{0.3cm}\noindent\rule{\linewidth}{0.5mm}\par\vspace{0.3cm}
{\huge\bfseries Arena\par}
\vspace{0.18cm}{\Large\color{aauRed}#1\par}
\vspace{0.15cm}{\normalsize\color{darkGray}#2\par}
\vspace{0.12cm}{\color{aauRed}\rule{\linewidth}{1.2pt}\par}
\vspace{0.35cm}{\normalsize\bfseries Prepared by\par}\vspace{0.15cm}
\begin{tabular}{@{}lr@{}}\textbf{Name}&\textbf{Student ID}\\\midrule
Khaled Alkhawaldeh&202410091\\Abdulrahman Alnaqbi&202340042\\Yousif Alhammadi&202310707\\Ahmed Alsamma&202320192\end{tabular}\par
\vspace{0.22cm}{\color{aauRed!55}\rule{0.6\textwidth}{1.2pt}\par}
\vspace{0.15cm}{\normalsize\textbf{Section:} 61D\par}
\vspace{0.15cm}{\normalsize\textbf{Submitted to:} Dr. Yazeed Ghadi\par}
\vspace{0.3cm}
\begin{tcolorbox}[colback=gray!5,colframe=black!70,title=\textbf{Project Summary},fonttitle=\bfseries,width=0.95\textwidth,boxrule=0.4mm,arc=3mm,left=7pt,right=7pt,top=4pt,bottom=4pt]
A Java desktop workspace for organizing campus competitions, coordinating teams, collecting submissions, evaluating work, and publishing results. Project, programming, and quiz scoring are separated from score-ranking, knockout, and league progression through object-oriented strategies. The application is developed for Al Ain University in the United Arab Emirates.
\end{tcolorbox}
\vfill{\normalsize Fall, 2026--2027\par}{\small Version 2.0 \;|\; \reportdate\par}
\end{titlepage}\pagenumbering{roman}\hypersetup{pageanchor=true}\tableofcontents\clearpage\pagenumbering{arabic}}

\renewcommand{\reportdate}{10 November 2026}
\begin{document}
\projecttitle{Design}{Architecture, object design, and format strategies}
\section{Introduction}
\subsection{Purpose of the System}
Arena translates campus competition processes into an executable object model. Its design places correctness in a Java application service and uses a Swing boundary to present those operations. Different formats share accounts, teams, registration, persistence, and publication while retaining distinct validation, evaluation, and progression policies.
\subsection{Design Goals}
\begin{itemize}
\item Make illegal operations fail before persistence: authorization, eligibility, stage readiness, and score bounds are explicit invariants.
\item Demonstrate object-oriented variation through separate scoring and tournament interfaces with four scoring policies and three tournament policies, rather than one generic score field.
\item Keep models immutable and persist complete, internally consistent updates.
\item Make the project buildable with a JDK alone and understandable during an OOAD demonstration.
\end{itemize}

\begin{umlpage}
\section{Software Architecture}
\subsection{Subsystem Decomposition}
\umlfigure{06_software_architecture}{Logical layers and their main responsibilities.}
\end{umlpage}
\textbf{Boundary: ui package.} ArenaApp owns the frame and navigation. AppDialogs owns forms; CompetitionViews builds cards, brackets, and leaderboards; UiComponents supplies shared controls; ArenaTheme defines the visual identity; WorkspacePanel tracks the viewport width while permitting vertical scrolling. It does not calculate official scores or directly mutate persisted records. UI work is scheduled on the Swing event-dispatch thread.

\textbf{Control: ArenaService.} Owns the active session and implements use cases. It validates role, ownership, phase, and domain constraints, delegates score policy to a strategy, and performs durable updates through Repository.

\textbf{Domain: Model, CompetitionRules, and TournamentRules.} Immutable records represent domain facts. CompetitionType selects a scoring policy; TournamentFormat selects fixture generation, advancement, completion, and match standings. CompetitionRules has ProjectRules, ProgrammingRules, QuizRules, and the legacy-named KnockoutRules no-rubric policy. TournamentRules has ScoreRankingRules, KnockoutTournamentRules, and RoundRobinRules. Phase is an enum with guarded forward transitions; this is an explicit state machine, not an implementation of the GoF State pattern.

\textbf{Infrastructure: Repository and Passwords.} Repository handles filtered serialization, file locking, and replacement. Passwords handles random salt generation and credential derivation. DemoData provides fictional fixtures only for an absent database.
\subsection{Hardware/Software Mapping}
One JVM process runs on a desktop with a graphical display, keyboard, mouse, a writable data directory, and JDK/JRE 17 or newer. The tested build used Windows and a portable JDK 21 with Java 17 compilation compatibility. No database server, web server, or network connection is required at runtime. Cross-platform Swing compatibility is an implementation goal; Linux/macOS runtime behavior has not been tested here.
\subsection{Persistent Data Management}
\code{Model.State} stores maps keyed by stable string IDs: users, teams, competitions, registrations, submissions, evaluations, and matches, plus an audit list and schema version. Records refer to one another by ID, avoiding recursive object graphs and accidental duplicate ownership. Insertion-ordered maps make registration order deterministic for bracket seeding.

Repository opens an exclusive operating-system lock on a companion file. It accepts schema versions 1 and 2, maps version-one formats to the new type/format fields in memory, writes schema two on a successful save, checks the file's 20 MB limit, and filters deserialization by application model classes and Java base classes, with depth, reference, and array limits. It does not accept arbitrary uploaded databases.

Each service mutation follows this transaction sequence:
\begin{enumerate}
\item Validate the request and serialize an in-memory rollback snapshot.
\item Apply the domain mutation and append a successful audit event.
\item Validate serialized size and readability, create a new temporary file, flush it, and replace the current file atomically where supported. A pre-existing temporary path is preserved and causes the update to fail.
\item If the operation or save fails, restore the in-memory snapshot and return an error.
\end{enumerate}
The repository is not a relational database. The tested version-one-to-two migration is provided; there is no general migration framework, write-ahead logging, encryption, backup rotation, or power-loss guarantee. The original application is not expected to read new quiz/league records. Preserve a closed-file backup before sharing a database with an older build. Copy the data file while the application is closed for a consistent manual backup. A corrupt existing file is preserved and startup stops.
\subsection{Access Control and Security}
\begin{tabularx}{\linewidth}{Ycccc}\toprule Operation & Student & Judge & Organizer & Admin\\\midrule
Browse / published results & Yes & Yes & Yes & Yes\\
Create/join a team & Yes & No & No & No\\
Register / submit & Captain & No & No & No\\
Create competition & No & No & Yes & Yes\\
Edit / advance & No & No & Owner & Any\\
Evaluate / decide match & No & Assigned & No & No\\
Create accounts & No & No & No & Yes\\
Read audit & No & No & Yes & Yes\\\bottomrule\end{tabularx}
Administrators manage competitions but do not impersonate judges. The service has a private session ID set by password verification; UI controls cannot supply an arbitrary actor for a mutation. Password records contain the iteration count, a random 16-byte salt, and a 256-bit PBKDF2-HMAC-SHA256 derivation using 120,000 iterations. Comparison uses \code{MessageDigest.isEqual}.

This is a local coursework security boundary, not internet-facing account security. There is no login throttling, reset email, institutional identity verification, or encrypted storage. All demo accounts use the documented initial password and should be changed if the dataset is used beyond a demonstration. URLs are validated as HTTP(S) and are not automatically opened or fetched. User-supplied names are displayed as literal text in labels, tables, and selection controls; intentional wrapped labels and tooltips escape markup.
\subsection{Boundary Conditions}
\begin{itemize}
\item First start: seed fictional accounts, teams, and competitions only if the data file does not exist.
\item Existing data: load without reseeding or resetting deadlines.
\item Duplicate process: fail with a lock error rather than permit concurrent writes.
\item Invalid action: show a message, preserve the form where possible, and keep the domain unchanged.
\item Closing: release the file lock. There is no unsaved-change prompt because every successful mutation is saved.
\item Incomplete judging or missing submissions: block the stage transition. Automatic forfeits and cancellation are future features.
\end{itemize}

\clearpage\section{Subsystem Services}
\subsection{User Management Subsystem}
\code{login(email,password)} clears the prior session before verification. Account creation requires an administrator and a unique email address.

\code{changePassword(old,new)} verifies the current credential and persists the new derivation. User directory views expose names, emails, and roles in the UI; credential hashes are never returned by account APIs. Login, current-user, and account-creation responses use the four-field UserSummary record. The full account directory is administrator-only; a separate organizer/admin judge picker returns only judge summaries.
\subsection{Competition Management Subsystem}
\code{createCompetition(...)} validates capacity, roster maximum, future deadline, and judge identities. \code{editCompetition(...)} is draft-only. \code{advance(cid)} applies guards before moving to the next enum value. Starting delegates fixture generation to TournamentRules inside the same transaction as the phase change. Unsupported type/format combinations fail before any mutation.
\subsection{Registration Subsystem}
\code{register(cid,tid)} checks the captain, phase, deadline, capacity, size, duplicate team, and overlapping members. \code{withdraw(cid,tid)} removes an entry only during registration. \code{joinTeam} and \code{leaveTeam} refuse changes to any registered roster. Students can belong to multiple teams globally but cannot enter the same competition twice.
\subsection{Submission and Evaluation Subsystem}
\code{submit(...)} replaces the one submission for a registered entry while Active and before its deadline. \code{evaluate(...)} replaces only the calling judge's review; other judges' reviews remain independent. Project scoring requires all assigned judges, while programming and quiz ranking permit one official score owner per team. Quiz uses a single bounded points field and does not require a submission.

\code{recordWinner(matchId,winnerId)} validates assignment and participant identity, stores the final winner, and creates the next round only once all current-round winners exist. The knockout strategy supports every count from 2 to 16 and records automatic first-round byes. \code{recordMatchScore(matchId,left,right)} records quiz knockout or league scores, rejects a knockout draw, and atomically saves the result and any generated next round. Round-robin fixtures are generated once and never duplicated after results.
\subsection{Results Management Subsystem}
\code{standings(cid)} enforces visibility, computes the selected format's scores, applies tie policies, and returns immutable Standing values. \code{standingsCsv(cid)} uses the same ordered values, quotes text cells, doubles embedded quotes, and neutralizes leading formula characters. Export does not maintain a second, potentially inconsistent ranking calculation.

\clearpage\section{Low-Level Design}
\subsection{Object Design Trade-offs}
\textbf{Composition over a user subclass hierarchy.} A User contains a Role enum because all roles share identity fields. Different permissions are checked by service operations. Artificial Student/Organizer subclasses would add little behavior in this implementation.

\textbf{Immutable records over mutable entity setters.} Team membership, judge sets, and competition fields are defensively copied. Changes create replacement records inside a service transaction. This reduces accidental state changes and makes rollback straightforward.

\textbf{Strategies plus explicit workflow branching.} CompetitionRules encapsulates scoring validation, score calculation, explanations, and submission requirements. TournamentRules separately encapsulates fixture generation, round advancement, completion, and match standings. The service retains explicit guards and score aggregation, including programming penalties. Extension requires the relevant enum, compatibility matrix, factory, UI explanation, and tests; this is a proportionate use of Strategy rather than a claim of modification-free extension.

\textbf{Snapshot persistence over SQL.} A local object snapshot avoids dependencies and setup. It is suitable for a small demonstration, but each write serializes the entire dataset. A future multi-user version should introduce a transactional database and server-side sessions.
\begin{umlpage}
\subsection{Final Object Design}
\umlfigure{07_strategy_class}{Independent scoring and tournament strategy interfaces and their concrete policies.}
\clearpage\subsubsection*{Design Class Diagram}
\umlfigure{08_design_class}{Complete source-checked overview of all 37 named production types. The vector figure can be zoomed; the following focused sheets provide readable class detail.}
\clearpage\subsubsection*{Application and UI Classes}
\umlfigure{08a_application_classes}{Application, UI, persistence, and support classes. Field and operation names match the Java source.}
\clearpage\subsubsection*{Records and Relationships}
\umlfigure{08b_record_classes}{Record fields and identifier-based relationships. Quiz evaluations attach to registered entries without requiring a submission.}
\end{umlpage}
\begin{umlpage}
\subsection{Packages and Dependencies}
\umlfigure{09_package_dependencies}{Main package dependencies. The UI also reads rule explanations, and ArenaApp.main constructs Repository at startup.}
\end{umlpage}
Production sources are under \code{src/main/java/edu/aau/arena}, with the serialized Model compatibility container under \code{src/main/java/edu/aau/ccms}; tests use \code{src/test/java/edu/aau/arena}. The application has no external runtime library. Maven uses JUnit Jupiter and Surefire to execute the existing lifecycle suites; the native scripts remain a JDK-only alternative. The ui, service, repository, rules, util, and model packages separate responsibilities. Existing serialized domain records intentionally retain their names under edu.aau.ccms.Model so version-one data stays readable; TC-69 verifies a pre-refactor fixture.
\subsection{Class Interfaces: Attributes, Methods, Responsibilities}
\small
\begin{longtable}{p{0.28\linewidth}p{0.64\linewidth}}\toprule Type & Selected real fields and operations\\\midrule\endhead
ArenaApp & service: ArenaService; main(args): void; showCompetition(cid): void\\
AppDialogs & app: ArenaApp; service: ArenaService; competitionDialog(c): void\\
CompetitionViews & app: ArenaApp; service: ArenaService; bracket(c), fixtures(c): JPanel\\
ArenaTheme & ACCENT, GREEN, PAPER: Color; actionButton(text, primary): JButton\\
UiComponents & grid(columns, rows): JTable; paragraph(text, background): JTextArea\\
WorkspacePanel & getScrollableTracksViewportWidth(): boolean\\
ArenaService & state: State; repository: Repository; session: String; clock: Clock; createCompetition(...): Competition; recordMatchScore(id, left, right): void; standings(cid): List<Standing>\\
Repository & file: Path; channel: FileChannel; lock: FileLock; load(): State; save(state): void; close(): void\\
DemoData & populate(state, now): void\\
Passwords & hash(password): String; verify(password, encoded): boolean\\
UserSummary & id, name, email: String; role: Role\\
CompetitionRules & validate(a, b, c): void; score(e): double; requiresAllJudges(): boolean\\
ProjectRules & score(e): double; requiresAllJudges(): boolean\\
ProgrammingRules & validate(a, b, c): void; score(e): double\\
QuizRules & maximum: int; validate(a, b, c): void\\
KnockoutRules & requiresSubmission(): boolean\\
TournamentRules & start(cid, teamIds): List<Match>; nextRound(cid, matches): List<Match>; matchesComplete(matches): boolean\\
ScoreRankingRules & start(cid, teams): List<Match>\\
Knockout\allowbreak TournamentRules & nextRound(cid, matches): List<Match>\\
RoundRobinRules & table(teams, matches, name):;     List<LeagueStanding>\\
Model & legacyType(format): CompetitionType; legacyTournament(format): TournamentFormat\\
State & schema: int; users, teams, competitions: Map; registrations, submissions: Map; evaluations, matches: Map; audit: List<Audit>; upgradeSchema(): void\\
User & id, name, email: String; role: Role; password: String\\
Team & id, name, captainId: String; members: List<String>\\
Competition & id, title, description: String; type: CompetitionType; tournamentFormat: TournamentFormat; phase: Phase; format: Format (legacy); capacity, teamSize, quizMaxPoints: int; organizerId: String; judgeIds: Set<String>; deadline: LocalDateTime; withPhase(next): Competition\\
Registration & id, competitionId, teamId: String\\
Submission & id, competitionId, teamId: String; title, description, url: String; submittedAt: LocalDateTime\\
Evaluation & id, competitionId, teamId, judgeId: String; a, b, c: int; feedback: String\\
Match & id, competitionId: String; round: int; leftTeamId, rightTeamId, winnerId: String; leftScore, rightScore: Integer; drawn: boolean; result(left, right): Match; bye(), decided(): boolean\\
Standing & rank: int; teamId, team: String; score: double; detail: String\\
LeagueStanding & rank, played, won, drawn, lost: int; scored, conceded, points: int; teamId, team: String; difference(): int\\
Audit & time: LocalDateTime; actor, action: String\\
Role & ADMIN, ORGANIZER, JUDGE, STUDENT\\
Format & PROJECT, PROGRAMMING, KNOCKOUT; Compatibility field for version-one data\\
CompetitionType & PROJECT, PROGRAMMING, QUIZ, MATCH; supports(format): boolean\\
TournamentFormat & SCORE\_BASED, KNOCKOUT, ROUND\_ROBIN\\
Phase & DRAFT, REGISTRATION, ACTIVE,; JUDGING, PUBLISHED, ARCHIVED\\\bottomrule\end{longtable}\normalsize
\subsection*{Invariants and Extension Example}
For every registration, the team and competition exist. A student's membership overlaps no other entry in that competition. For every project/programming scored entry, a submission exists. Quiz ranking entries need no submission. For each project entry there is at most one evaluation per judge. For programming and quiz ranking there is at most one official evaluation per entry. Every match winner belongs to its match. A published competition cannot accept updated scores.

To add a new scored activity, define a CompetitionType and scoring policy, update its allowed tournament combinations and factory, and add service/UI tests. To add another tournament system, implement TournamentRules and update the TournamentFormat factory and compatibility checks. Round robin is already implemented independently of quiz scoring.

\clearpage\section{Glossary \& References}
\textbf{Facade:} a single application-level entry point for use cases. \textbf{Strategy:} an interchangeable rule implementation. \textbf{Record:} an immutable Java data carrier. \textbf{Guard:} a precondition for a transition. \textbf{Rollback snapshot:} a serialized copy used to restore memory after a failed write. \textbf{EDT:} Swing event-dispatch thread.
\begingroup\raggedright\begin{enumerate}
\item Object Management Group. \emph{Unified Modeling Language, version 2.5.1}. \url{https://www.omg.org/spec/UML/2.5.1}. Accessed 29 October 2026. Used as the notation reference for the UML models.
\item Oracle. \emph{Java SE 17: javax.swing package}. \url{https://docs.oracle.com/en/java/javase/17/docs/api/java.desktop/javax/swing/package-summary.html}. Accessed 29 October 2026. Used for desktop components and event-dispatch-thread conventions.
\item Oracle. \emph{Java SE 17: ObjectInputFilter}. \url{https://docs.oracle.com/en/java/javase/17/docs/api/java.base/java/io/ObjectInputFilter.html}. Accessed 29 October 2026. Used for bounded, filtered repository deserialization.
\item Project materials. Task outlines in \code{requirements/system-analysis-requirements.txt}, \code{requirements/system-design-requirements.txt}, \code{requirements/implementation-report-requirements.txt}, \code{requirements/testing-requirements.txt}, and \code{requirements/sample-data-requirements.txt}; instructor feedback supplied by the team. These define the coursework scope, not an independent evaluation of the implementation.
\end{enumerate}\endgroup

\end{document}
